\chapter{Méthodologie et organisation}

\section{Gestion de projet avec GitHub}

Dans cette section, je vais vous présenter mon approche de gestion de projet entièrement centralisée sur GitHub. Vous verrez ma maîtrise des outils GitHub pour la planification, le suivi et la réalisation de mon projet.
\newline

\begin{for_you}{Mon approche GitHub :}
    \textit{J'utilise GitHub pour stocker et partager mes projets,
    GitHub Projects pour gérer la partie organisation et versionnement, et GitHub Actions pour automatiser certaines tâches.}
\end{for_you}

\subsection{User Stories et estimation de temps}

Dans cette sous-section, vous trouverez le détail de mes User Stories avec des estimations de temps réalistes. Chaque User Story est liée à des commits et des milestones GitHub pour un suivi précis de l'avancement.
\newline

\begin{for_you}{Kanban :}
\begin{verbatim}
https://github.com/users/Theoden-bernard/projects/2/views/2
\end{verbatim}

\textit{}

\textbf{Colonnes du tableau Kanban :}
\begin{itemize}
    \item \textbf{Backlog :} Fonctionnalités à développer
    \item \textbf{Ready :} Tâches prêtes pour le sprint
    \item \textbf{In Progress :} Tâches en cours (WIP limit : 3)
    \item \textbf{In Review :} Code en attente de validation
    \item \textbf{Done :} Fonctionnalités livrées
\end{itemize}
\end{for_you}

\section{Versioning GitHub et conventions}

Le versioning GitHub suit le modèle Git Flow avec des branches spécialisées pour chaque type de développement. Les conventions de nommage et de commit facilitent la traçabilité et la collaboration. Les Pull Requests permettent la revue de code systématique et la validation des fonctionnalités avant intégration.

Les conventions établies couvrent le nommage des branches, le format des messages de commit, et les templates de Pull Request. Cette standardisation améliore la qualité du code et accélère l'onboarding de nouveaux développeurs.



\begin{for_you}{Norme GitHub}
\textbf{Schéma Git Flow :}
\begin{verbatim}
main -----------------------------------------------
  |
  +-- develop --------------------------------------
       |
       +-- feature/user-authentication
       +-- feature/project-management
       +-- hotfix/critical-bug-fix
\end{verbatim}

\textbf{Conventions de commit :}
\begin{lstlisting}
FEAT: add user authentication system
FIX: resolve login validation issue
DOCS: update documentation
TEST: add unit tests for user service
REFACTOR: improve code structure
\end{lstlisting}
\end{for_you}

\section{Planification et outils de suivi}

La planification combine une roadmap GitHub pour la vision macro et GitHub Projects pour le suivi opérationnel. La roadmap GitHub visualise les dépendances et les jalons critiques, tandis que le Kanban GitHub Projects offre une vue détaillée des tâches en cours. Cette approche duale optimise la coordination entre la planification stratégique et l'exécution tactique.

La roadmap GitHub permet de communiquer la vision produit et les priorités à long terme. Les milestones et les dépendances facilitent la coordination entre les différentes équipes et la gestion des risques de planning.
\newline

\begin{for_you}{Roadmap GitHub :}
\begin{verbatim}
Phase 1: Conception (4 semaines)
+-- Analyse des besoins (1 semaine)
+-- Conception technique (2 semaines)
+-- Validation architecture (1 semaine)

Phase 2: Développement MVP (8 semaines)
+-- Database (L)
+-- Authentification (S)
+-- Frontend Phoenix (S)
+-- Badge (L)
\end{verbatim}

\textbf{\newline Configuration GitHub Projects :}
\begin{itemize}
    \item \textbf{Colonnes :} Backlog, To Do, In Progress, Review, Done
    \item \textbf{WIP Limits :} 3 tâches max en cours par développeur
    \item \textbf{Policies :} PR obligatoire pour merge en develop
    \item \textbf{Automation :} Mise à jour automatique des statuts
\end{itemize}
\end{for_you}

\newpage
\section{Estimation de temps et planification}

Dans cette section, je vais vous présenter mes estimations de temps pour chaque fonctionnalité et expliquer comment j'ai planifié mon projet.
Nous verrons une approche réaliste et méthodique de la gestion du temps.
\newline
L'analyse des temps permet de valider la faisabilité du projet et d'optimiser la planification selon les contraintes disponibles. Cette approche pragmatique démontre votre capacité à prendre en compte les contraintes temporelles dans les décisions techniques.
\newline

\begin{for_you}{Estimation de temps par fonctionnalité :}
    \begin{longtable}{p{3cm}p{3.2cm}p{0.9cm}p{0.6cm}p{1.6cm}p{1.7cm}}
    \toprule
    \textbf{Fonctionnalité} & \textbf{Description} & \textbf{Unité} & \textbf{Qty} & \textbf{Temps estimé} & \textbf{Total} \\
        \midrule
        Authentification & Login/Register & jours & 3 & 2 & 6 \\
        Gestion projets & CRUD projets & jours & 5 & 3 & 15 \\
        Tableaux de bord & Dashboard & jours & 2 & 4 & 8 \\
        Tests & Tests unitaires & jours & 10 & 1 & 10 \\
        Déploiement & CI/CD & jours & 3 & 2 & 6 \\
        \midrule
    \multicolumn{5}{r}{\textbf{Total estimé}} & \textbf{45 jours} \\
    \bottomrule
    \end{longtable}
\end{for_you}

\section{Liens utiles}

\begin{itemize}
    \item GitHub Flow/PRs: \url{https://docs.github.com/pull-requests}
    \item Git Flow: \url{https://bit.ly/gitflow-atlassian}
    \item GitHub Projects: \url{https://bit.ly/github-projects}
    \item GitHub Roadmap: \url{https://bit.ly/github-roadmap}
    \item GitHub Milestones: \url{https://bit.ly/github-milestones}
    \item User Stories: \url{https://www.mountaingoatsoftware.com/agile/user-stories}
    \item Estimation de temps: \url{https://bit.ly/time-estimation}
\end{itemize}